\textbf{Phần 1. Trắc nghiệm lựa chọn}
\setcounterpageref{ex}{0}

\begin{ex} %[4D1V05]
Một đội xe chở cát gồm 3 tổ. Tổ 1 chở được $2400\text{ kg}$ cát. Tổ 2 chở được nhiều hơn tổ 1 là $600\text{ kg}$ cát. Tổ 3 chở được số cát gấp đôi hiệu số cát của tổ 2 và tổ 1. Cả ba tổ chở được tất cả số ki-lô-gam cát là
\choice
{$6000\text{ kg}$}
{\True $6600\text{ kg}$}
{$7200\text{ kg}$}
{$5400\text{ kg}$}
\loigiai{
Tổ 2 chở được số cát là:
$$2400 + 600 = 3000\text{ (kg)}$$
Hiệu số cát của tổ 2 và tổ 1 là $600\text{ kg}$. Tổ 3 chở được số cát là:
$$600 \times 2 = 1200\text{ (kg)}$$
Cả ba tổ chở được tất cả số ki-lô-gam cát là:
$$2400 + 3000 + 1200 = 6600\text{ (kg)}$$
\textbf{Đáp án của câu hỏi là B}
}
\end{ex}

\begin{ex} %[4D1V05]
Một kho gạo có $16\,500\text{ kg}$ gạo. Ngày thứ nhất kho xuất $3500\text{ kg}$ gạo. Ngày thứ hai kho xuất số gạo gấp đôi ngày thứ nhất. Số gạo còn lại trong kho chia đều vào $8$ xe tải để chở đến các siêu thị. Mỗi xe tải chở số ki-lô-gam gạo là
\choice
{$800\text{ kg}$}
{$650\text{ kg}$}
{$700\text{ kg}$}
{\True $750\text{ kg}$}
\loigiai{
Ngày thứ hai kho đã xuất số gạo là:
$$3500 \times 2 = 7000\text{ (kg)}$$
Cả hai ngày kho đã xuất số gạo là:
$$3500 + 7000 = 10\,500\text{ (kg)}$$
Số gạo còn lại trong kho là:
$$16\,500 - 10\,500 = 6000\text{ (kg)}$$
Mỗi xe tải chở số ki-lô-gam gạo là:
$$6000 : 8 = 750\text{ (kg)}$$
\textbf{Đáp án của câu hỏi là D}
}
\end{ex}

\begin{ex} %[4D1V05]
Mẹ mua $4$ hộp bánh dẻo hết $120\,000$ đồng và $3$ hộp bánh nướng cùng loại, biết giá tiền mỗi hộp bánh nướng đắt hơn mỗi hộp bánh dẻo là $15\,000$ đồng. Mẹ đưa cho người bán hàng tờ tiền $500\,000$ đồng. Người bán hàng phải trả lại mẹ số tiền là
\choice
{\True $245\,000$ đồng}
{$255\,000$ đồng}
{$235\,000$ đồng}
{$265\,000$ đồng}
\loigiai{
Giá tiền của một hộp bánh dẻo là:
$$120\,000 : 4 = 30\,000\text{ (đồng)}$$
Giá tiền của một hộp bánh nướng là:
$$30\,000 + 15\,000 = 45\,000\text{ (đồng)}$$
Số tiền mẹ mua 3 hộp bánh nướng là:
$$45\,000 \times 3 = 135\,000\text{ (đồng)}$$
Tổng số tiền mẹ đã mua bánh là:
$$120\,000 + 135\,000 = 255\,000\text{ (đồng)}$$
Người bán hàng phải trả lại mẹ số tiền là:
$$500\,000 - 255\,000 = 245\,000\text{ (đồng)}$$
\textbf{Đáp án của câu hỏi là A}
}
\end{ex}

\begin{ex} %[4D1V05]
Một trang trại nuôi gà ở 3 chuồng. Chuồng A có $1250$ con gà. Chuồng B có số gà gấp 3 lần chuồng A. Chuồng C có số gà bằng trung bình cộng số gà của chuồng A và chuồng B. Cả ba chuồng có tất cả số con gà là
\choice
{$6500$ con}
{$7000$ con}
{\True $7500$ con}
{$8000$ con}
\loigiai{
Chuồng B có số gà là:
$$1250 \times 3 = 3750\text{ (con)}$$
Chuồng C có số gà là:
$$(1250 + 3750) : 2 = 2500\text{ (con)}$$
Cả ba chuồng có tất cả số gà là:
$$1250 + 3750 + 2500 = 7500\text{ (con)}$$
\textbf{Đáp án của câu hỏi là C}
}
\end{ex}

\begin{ex} %[4D1V05]
Một đội công nhân cần sửa đoạn đường dài $4800\text{ m}$. Tuần thứ nhất đội sửa được $\frac{1}{4}$ đoạn đường đó. Tuần thứ hai đội sửa được số mét đường gấp đôi tuần thứ nhất. Sau hai tuần, đội đó còn phải sửa số mét đường nữa để hoàn thành là
\choice
{$1000\text{ m}$}
{\True $1200\text{ m}$}
{$1500\text{ m}$}
{$2400\text{ m}$}
\loigiai{
Tuần thứ nhất đội sửa được:
$$4800 : 4 = 1200\text{ (m)}$$
Tuần thứ hai đội sửa được:
$$1200 \times 2 = 2400\text{ (m)}$$
Cả hai tuần đội đã sửa được:
$$1200 + 2400 = 3600\text{ (m)}$$
Số mét đường đội còn phải sửa tiếp là:
$$4800 - 3600 = 1200\text{ (m)}$$
\textbf{Đáp án của câu hỏi là B}
}
\end{ex}

\begin{ex} %[4D1V05]
Bác Tư thu hoạch được $1800\text{ kg}$ cà phê. Bác Năm thu hoạch được ít hơn bác Tư $300\text{ kg}$. Bác Sáu thu hoạch được nhiều hơn mức trung bình cộng của cả ba bác là $100\text{ kg}$. Bác Sáu thu hoạch được số ki-lô-gam cà phê là
\choice
{$1600\text{ kg}$}
{$1700\text{ kg}$}
{$1750\text{ kg}$}
{\True $1800\text{ kg}$}
\loigiai{
Bác Năm thu hoạch được số cà phê là:
$$1800 - 300 = 1500\text{ (kg)}$$
Tổng số cà phê của bác Tư và bác Năm là:
$$1800 + 1500 = 3300\text{ (kg)}$$
Mức trung bình cộng của cả ba bác là:
$$(3300 + 100) : 2 = 1700\text{ (kg)}$$
Bác Sáu thu hoạch được số cà phê là:
$$1700 + 100 = 1800\text{ (kg)}$$
\textbf{Đáp án của câu hỏi là D}
}
\end{ex}

\begin{ex} %[4D1V05]
Ba xe chở khách: Xe thứ nhất chở $45$ người. Nếu chuyển $5$ người từ xe thứ hai sang xe thứ nhất thì số khách ở hai xe bằng nhau. Xe thứ ba chở số khách bằng một nửa tổng số khách của hai xe đầu lúc ban đầu. Cả ba xe chở tất cả số người là
\choice
{\True $150$ người}
{$140$ người}
{$160$ người}
{$145$ người}
\loigiai{
Khi chuyển 5 người từ xe thứ hai sang xe thứ nhất thì hai xe bằng nhau, tức là ban đầu xe thứ hai nhiều hơn xe thứ nhất:
$$5 \times 2 = 10\text{ (người)}$$
Số khách ban đầu ở xe thứ hai là:
$$45 + 10 = 55\text{ (người)}$$
Tổng số khách ở hai xe đầu là:
$$45 + 55 = 100\text{ (người)}$$
Số khách ở xe thứ ba là:
$$100 : 2 = 50\text{ (người)}$$
Cả ba xe chở được tất cả:
$$100 + 50 = 150\text{ (người)}$$
\textbf{Đáp án của câu hỏi là A}
}
\end{ex}

\begin{ex} %[4D1V05]
Một cửa hàng có $5$ hộp bút mực, mỗi hộp chứa $12$ chiếc và $4$ hộp bút bi, mỗi hộp chứa $20$ chiếc. Cửa hàng đã bán được $\frac{1}{2}$ số bút mực và $\frac{1}{4}$ số bút bi. Cửa hàng còn lại tất cả số chiếc bút là
\choice
{$80$ chiếc}
{$85$ chiếc}
{\True $90$ chiếc}
{$100$ chiếc}
\loigiai{
Tổng số bút mực cửa hàng có là:
$$5 \times 12 = 60\text{ (chiếc)}$$
Số bút mực còn lại là:
$$60 - 60 : 2 = 30\text{ (chiếc)}$$
Tổng số bút bi cửa hàng có là:
$$4 \times 20 = 80\text{ (chiếc)}$$
Số bút bi còn lại là:
$$80 - 80 : 4 = 60\text{ (chiếc)}$$
Tổng số chiếc bút còn lại là:
$$30 + 60 = 90\text{ (chiếc)}$$
\textbf{Đáp án của câu hỏi là C}
}
\end{ex}

\begin{ex} %[4D1V05]
Người ta có $100\,000$ đồng đi chợ mua $2\text{ kg}$ khoai tây với giá $15\,000$ đồng/kg và $3\text{ kg}$ cà chua với giá $18\,000$ đồng/kg. Số tiền còn lại dùng để mua các mớ rau cải với giá $8000$ đồng một mớ. Người đó mua được tối đa số mớ rau cải là
\choice
{$1$ mớ}
{\True $2$ mớ}
{$3$ mớ}
{$4$ mớ}
\loigiai{
Số tiền mua khoai tây là:
$$2 \times 15\,000 = 30\,000\text{ (đồng)}$$
Số tiền mua cà chua là:
$$3 \times 18\,000 = 54\,000\text{ (đồng)}$$
Số tiền còn lại sau khi mua khoai tây và cà chua là:
$$100\,000 - (30\,000 + 54\,000) = 16\,000\text{ (đồng)}$$
Số mớ rau cải mua được là:
$$16\,000 : 8000 = 2\text{ (mớ)}$$
\textbf{Đáp án của câu hỏi là B}
}
\end{ex}

\begin{ex} %[4D1V05]
Có 3 thùng chứa nước mắm: Thùng thứ nhất chứa $45\text{ l}$. Thùng thứ hai nếu bớt đi $5\text{ l}$ thì gấp đôi thùng thứ nhất. Thùng thứ ba chứa số lít nước mắm bằng một nửa tổng số nước mắm ở hai thùng đầu. Cả ba thùng chứa số lít nước mắm là
\choice
{$180\text{ l}$}
{$195\text{ l}$}
{$200\text{ l}$}
{\True $210\text{ l}$}
\loigiai{
Thùng thứ hai chứa số lít nước mắm là:
$$45 \times 2 + 5 = 95\text{ (l)}$$
Tổng số nước mắm ở hai thùng đầu là:
$$45 + 95 = 140\text{ (l)}$$
Thùng thứ ba chứa số lít nước mắm là:
$$140 : 2 = 70\text{ (l)}$$
Cả ba thùng chứa tất cả:
$$140 + 70 = 210\text{ (l)}$$
\textbf{Đáp án của câu hỏi là D}
}
\end{ex}

\begin{ex} %[4D1V05]
Một nông trường nuôi gà và vịt, số con gà nhiều hơn số con vịt là $120$ con. Sau khi người ta bán đi $30$ con gà và mua thêm $10$ con vịt thì lúc này số con gà nhiều hơn số con vịt là
\choice
{\True $80$ con}
{$100$ con}
{$140$ con}
{$85$ con}
\loigiai{
Khi bán đi 30 con gà, hiệu số gà và vịt giảm đi 30 con:
$$120 - 30 = 90\text{ (con)}$$
Khi mua thêm 10 con vịt, hiệu số gà và vịt giảm tiếp 10 con:
$$90 - 10 = 80\text{ (con)}$$
Vậy lúc này số gà nhiều hơn số vịt $80$ con.\\
\textbf{Đáp án của câu hỏi là A}
}
\end{ex}

\begin{ex} %[4D1V05]
Hai bạn Hùng và Dũng có tất cả $60$ viên bi. Nếu Hùng cho Dũng $6$ viên bi thì số bi của Dũng gấp đôi số bi của Hùng lúc đó. Ban đầu Hùng có số viên bi là
\choice
{$20$ viên}
{$24$ viên}
{\True $26$ viên}
{$34$ viên}
\loigiai{
Khi Hùng cho Dũng 6 viên bi thì tổng số bi của hai bạn không thay đổi và vẫn là 60 viên bi.\\
Khi đó số bi của Hùng là 1 phần thì số bi của Dũng là 2 phần như thế.\\
Tổng số phần bằng nhau là:
$$1 + 2 = 3\text{ (phần)}$$
Số bi của Hùng sau khi đã cho Dũng là:
$$60 : 3 = 20\text{ (viên)}$$
Số bi ban đầu của Hùng là:
$$20 + 6 = 26\text{ (viên)}$$
\textbf{Đáp án của câu hỏi là C}
}
\end{ex}

\textbf{Phần 2. Trắc nghiệm đúng sai}
\setcounterpageref{ex}{0}

\begin{ex} %[4D1V05]
Một cửa hàng hoa quả nhập về $150\text{ kg}$ cam. Ngày thứ nhất cửa hàng bán được $40\text{ kg}$ cam. Ngày thứ hai bán được số cam gấp đôi ngày thứ nhất. Số cam còn lại đóng đều vào các túi, mỗi túi $5\text{ kg}$. Xét tính đúng sai của các khẳng định sau:
\choiceTF
{\True Ngày thứ hai cửa hàng bán được $80\text{ kg}$ cam.}
{Tổng số cam đã bán sau hai ngày là $110\text{ kg}$.}
{\True Số cam còn lại chưa bán là $30\text{ kg}$.}
{\True Đóng được tất cả đúng $6$ túi cam từ số cam còn lại.}
\loigiai{
\begin{itemize}
    \item Số cam bán ngày thứ hai: $40 \times 2 = 80\text{ (kg)}$. Khẳng định a Đúng.
    \item Tổng số cam đã bán: $40 + 80 = 120\text{ (kg)}$ (không phải $110\text{ kg}$). Khẳng định b Sai.
    \item Số cam còn lại: $150 - 120 = 30\text{ (kg)}$. Khẳng định c Đúng.
    \item Số túi cam đóng được: $30 : 5 = 6\text{ (túi)}$. Khẳng định d Đúng.
\end{itemize}
\begin{itemize}
    \item a) \textbf{Đúng}.
    \item b) \textbf{Sai}.
    \item c) \textbf{Đúng}.
    \item d) \textbf{Đúng}.
\end{itemize}
}
\end{ex}

\begin{ex} %[4D1V05]
Ba bạn An, Bình, Cường cùng góp tiền mua một quả bóng đá giá $180\,000$ đồng. An góp $\frac{1}{3}$ giá tiền quả bóng. Bình góp nhiều hơn An $20\,000$ đồng. Số tiền còn lại do Cường góp. Xét tính đúng sai của các khẳng định sau:
\choiceTF
{\True Số tiền An đã góp là $60\,000$ đồng.}
{\True Số tiền Bình đã góp là $80\,000$ đồng.}
{Số tiền Cường góp nhiều hơn số tiền An góp.}
{\True Số tiền Cường góp bằng một nửa số tiền Bình góp.}
\loigiai{
\begin{itemize}
    \item Số tiền An góp: $180\,000 : 3 = 60\,000\text{ (đồng)}$. Khẳng định a Đúng.
    \item Số tiền Bình góp: $60\,000 + 20\,000 = 80\,000\text{ (đồng)}$. Khẳng định b Đúng.
    \item Cả An và Bình góp: $60\,000 + 80\,000 = 140\,000\text{ (đồng)}$.\\
    Số tiền Cường góp: $180\,000 - 140\,000 = 40\,000\text{ (đồng)}$.\\
    Vì $40\,000 < 60\,000$ nên Cường góp ít hơn An. Khẳng định c Sai.
    \item Ta có $40\,000 = 80\,000 : 2$ nên Cường góp bằng một nửa Bình. Khẳng định d Đúng.
\end{itemize}
\begin{itemize}
    \item a) \textbf{Đúng}.
    \item b) \textbf{Đúng}.
    \item c) \textbf{Sai}.
    \item d) \textbf{Đúng}.
\end{itemize}
}
\end{ex}

\textbf{Phần 3. Trả lời ngắn}
\setcounterpageref{ex}{0}

\begin{ex} %[4D1V05]
Một kho có $4500$ bao xi măng. Ngày thứ nhất xuất đi $1200$ bao xi măng, ngày thứ hai xuất đi số bao gấp đôi ngày thứ nhất. Hỏi trong kho còn lại bao nhiêu bao xi măng?
\shortans{900}
\loigiai{
\textbf{Đáp số:} $900$\\
Số bao xi măng xuất trong ngày thứ hai là:
$$1200 \times 2 = 2400\text{ (bao)}$$
Cả hai ngày đã xuất số bao xi măng là:
$$1200 + 2400 = 3600\text{ (bao)}$$
Số bao xi măng còn lại trong kho là:
$$4500 - 3600 = 900\text{ (bao)}$$
}
\end{ex}

\begin{ex} %[4D1V05]
Ba bạn Mai, Hoa, Lan gấp được tất cả $135$ ngôi sao giấy. Mai gấp được $35$ ngôi sao. Hoa gấp được nhiều hơn Mai $10$ ngôi sao. Hỏi Lan gấp được bao nhiêu ngôi sao?
\shortans{55}
\loigiai{
\textbf{Đáp số:} $55$\\
Số ngôi sao Hoa gấp được là:
$$35 + 10 = 45\text{ (ngôi sao)}$$
Tổng số ngôi sao Mai và Hoa gấp được là:
$$35 + 45 = 80\text{ (ngôi sao)}$$
Số ngôi sao Lan gấp được là:
$$135 - 80 = 55\text{ (ngôi sao)}$$
}
\end{ex}

\begin{ex} %[4D1V05]
Một phân xưởng nhận về $3$ cuộn vải, mỗi cuộn dài $120\text{ m}$. Phân xưởng dùng số vải đó để may đồng phục học sinh, mỗi bộ hết $3\text{ m}$ vải. Sau khi may xong, phân xưởng còn thừa lại $30\text{ m}$ vải. Hỏi phân xưởng đã may được bao nhiêu bộ đồng phục?
\shortans{110}
\loigiai{
\textbf{Đáp số:} $110$\\
Tổng số mét vải trong 3 cuộn là:
$$3 \times 120 = 360\text{ (m)}$$
Số mét vải đã dùng để may đồng phục là:
$$360 - 30 = 330\text{ (m)}$$
Số bộ đồng phục đã may được là:
$$330 : 3 = 110\text{ (bộ)}$$
}
\end{ex}

\begin{ex} %[4D1V05]
Một thửa ruộng hình chữ nhật có chu vi $140\text{ m}$, chiều dài hơn chiều rộng $10\text{ m}$. Cứ mỗi $10\text{ m}^2$ thửa ruộng đó thu hoạch được $6\text{ kg}$ thóc. Hỏi cả thửa ruộng thu hoạch được bao nhiêu ki-lô-gam thóc?
\shortans{720}
\loigiai{
\textbf{Đáp số:} $720$\\
Nửa chu vi của thửa ruộng là:
$$140 : 2 = 70\text{ (m)}$$
Chiều dài thửa ruộng là:
$$(70 + 10) : 2 = 40\text{ (m)}$$
Chiều rộng thửa ruộng là:
$$40 - 10 = 30\text{ (m)}$$
Diện tích thửa ruộng là:
$$40 \times 30 = 1200\text{ (m}^2\text{)}$$
Khối lượng thóc thu hoạch được là:
$$1200 : 10 \times 6 = 720\text{ (kg)}$$
}
\end{ex}

\begin{ex} %[4D1V05]
Có ba can đựng dầu. Can thứ nhất đựng $18\text{ l}$ dầu. Can thứ hai đựng lượng dầu gấp đôi can thứ nhất. Can thứ ba đựng lượng dầu bằng trung bình cộng của cả ba can. Hỏi can thứ ba đựng bao nhiêu lít dầu?
\shortans{27}
\loigiai{
\textbf{Đáp số:} $27$\\
Can thứ hai đựng số lít dầu là:
$$18 \times 2 = 36\text{ (l)}$$
Vì can thứ ba đựng lượng dầu bằng trung bình cộng của cả ba can nên lượng dầu ở can thứ ba chính bằng trung bình cộng lượng dầu của hai can đầu.\\
Can thứ ba đựng số lít dầu là:
$$(18 + 36) : 2 = 27\text{ (l)}$$
}
\end{ex}

\begin{ex} %[4D1V05]
An có một số tiền. An mua sách hết một nửa số tiền và thêm $5000$ đồng. Sau đó An mua vở hết một nửa số tiền còn lại và thêm $5000$ đồng thì An còn lại đúng $10\,000$ đồng. Hỏi ban đầu An có bao nhiêu nghìn đồng?
\shortans{70}
\loigiai{
\textbf{Đáp số:} $70$\\
Số tiền của An trước khi mua vở là:
$$(10\,000 + 5000) \times 2 = 30\,000\text{ (đồng)}$$
Số tiền ban đầu của An là:
$$(30\,000 + 5000) \times 2 = 70\,000\text{ (đồng)}$$
Vậy ban đầu An có $70$ nghìn đồng.
}
\end{ex}
